\section{数据来源与证据层次}
\label{sec:01_04}

数据能够支持什么结论，取决于它如何产生、经过哪些转换，以及最终核验了什么对象。一个便于手算的例子可以解释约束，却不能代表实际区域的运行难度；公开拓扑可以提供可追溯的关联结构，却不一定包含巡检所需的运动、成像和能量参数；求解器报告了目标值，也不等于已经取得可执行的巡检方案。本书据此区分输入来源、实验计数和结果验收三个方面，分别说明模型的依据与结论的适用范围。

\subsection{输入来源及其作用}

本书使用的材料包括一般数学定义、人工教学例、受控合成输入、公开拓扑以及真实地理参考。求解与验收见证则记录算法在这些输入上实际产生了什么结果。它们并不是一条可以相互替代的证据等级链：输入来源回答模型建立在什么材料之上，见证核验回答某个具体输出是否满足指定要求。表\ref{s0104:evidence_sources}列出各类材料的作用。

\begin{table}[htbp]
 \centering
 \small
 \caption{数据来源、核验对象及其证据边界}
 \label{s0104:evidence_sources}
 \begin{tabularx}{\linewidth}{@{}p{0.15\linewidth}XX@{}}
  \toprule
  材料类别 & 提供的信息与用途 & 不能据此直接得到的结论\\
  \midrule
  一般定义 & 明确集合、关联、变量域和约束语义，建立可推导的数学关系 & 定义本身不证明某种算法在实际输入上的性能\\
  人工教学例 & 用有限且可逐项检查的参数解释选择、冲突和目标计算 & 小例子的最优性或算法差异不能直接外推到大规模输入\\
  受控合成输入 & 按给定生成规则构造关联结构，用于考察指定因素变化下的算法表现 & 生成规则不等于实际巡检任务的统计分布\\
  公开拓扑 & 提供可追溯的原始关联结构，并保留转换和赋权规则 & 使用论文输入不等于复现了该论文的全部算法\\
  真实地理参考 & 为区域位置、道路等空间关系提供依据，与目标模板和设备参考参数共同构成数值场景 & 地理位置真实不等于几何、设备和运行参数均经现场标定\\
  求解与验收见证 & 保存具体决策并检查容量、目标值或执行约束，区分可行性、覆盖程度与停止状态 & 单个可行见证不证明全局最优，数值验收不等于现场试飞\\
  \bottomrule
 \end{tabularx}
\end{table}

一般定义不依赖某一批实验。例如，候选是否共同占用某项资源，是关联关系的定义；某算法在这些关联关系上取得多大收益，则是需要运行记录支持的经验事实。人工教学例连接这两类内容：先固定所有参数，再把公式逐项代入，使读者能够独立核查一个选择为什么可行、目标值怎样得到。受控合成输入进一步扩大规模或改变结构，但其结论仍以生成规则和实际覆盖的条件为边界。

公开输入需要区分原始拓扑、转换后的优化问题和实际执行的方法。本书复用的六个最大独立集图来自\href{https://doi.org/10.1038/s42256-025-01141-4}{\textit{Reusability report: A distributed strategy for solving combinatorial optimization problems with hypergraph neural networks}}所对应的公开输入，经等价转换用于组合选择实验；HypOp神经网络求解器没有执行。因此，这六个输入提供的是问题结构来源，不能被写成六次神经方法实验，更不能据此比较神经方法与本书所运行方法的性能。

\href{https://doi.org/10.7155/jgaa.v30i1.3166}{\textit{Engineering Hypergraph b-Matching Algorithms}}提供了七个原始拓扑及相应的\texttt{random100}条件。相关实验既包含局部的Greedy与ILS诊断，也包含原作者完整流程的运行；二者的算法范围不同，必须分别报告。\href{https://doi.org/10.4230/LIPIcs.ESA.2025.79}{\textit{Semi-Streaming Algorithms for Hypergraph Matching}}对应的输入则保留了完整的\texttt{superblue19}拓扑，并依照公开源码重建LHG权重。具体地，对非空有限候选集合$E$，每个$e\in E$表示一个候选及其所含资源，令$r=\max_{e\in E}|e|$为最大资源数，则该赋权规则为
\begin{equation}
 w_e=r+1-|e|,\qquad e\in E.
 \label{s0104:public_weight}
\end{equation}
这里，$|e|$和$r$均为无量纲整数，$w_e$为无量纲收益分值，满足$w_e\geq1$。这一权重来自明确的重建规则，而不是现场测得的任务收益。保留拓扑来源和赋权规则，才能说明实验究竟复用了公开材料的哪一部分。

\subsection{实验条件、输入实例与运行次数}

实验条件、输入实例和运行次数具有不同含义。实验条件是清单中用于组织比较的一项输入及参数配置；输入实例是落实了资源、候选和权重等内容的具体优化对象；一次运行则是在给定条件下，使用某种方法及其随机种子、停止规则等配置执行一次计算。同一个输入可以对应多种条件，同一个条件也可以运行多个方法或多个种子。因此，条件数不能直接称为独立数据集数，重复运行数也不能直接称为独立实例数。

现有组合实验清单包含初始条件与压力条件两部分。分别记其有限索引集合为$\mathcal C_{\mathrm{init}}$和$\mathcal C_{\mathrm{stress}}$；索引仅用于标识清单条目，不具有物理单位。按当前清单口径，两部分不重复计数，其组成满足
\begin{subequations}
 \label{s0104:condition_counts}
 \begin{align}
 |\mathcal C_{\mathrm{init}}|
 &=1+18+6+7+1=33,
 \label{s0104:initial_counts}\\
 |\mathcal C_{\mathrm{stress}}|
 &=39+20=59,
 \label{s0104:stress_counts}\\
 |\mathcal C_{\mathrm{init}}\cup\mathcal C_{\mathrm{stress}}|
 &=33+59=92.
 \label{s0104:total_counts}
 \end{align}
\end{subequations}
式\eqref{s0104:initial_counts}中的五项依次对应一个人工教学例、十八个初始合成条件、六个HypOp来源的最大独立集输入条件、七个JGAA来源条件和一个ESA来源条件。式\eqref{s0104:stress_counts}中的两项分别为三十九个受控压力条件和二十个公开输入派生压力条件。派生条件可以用于检验结构或参数变化，但并不因此获得新的独立来源。

五十九个压力条件各有四种方法的运行记录，因而压力实验共有
\begin{equation}
 59\times4=236
 \label{s0104:stress_runs}
\end{equation}
次运行。独立算术核验记录确认了九十二项条件的输入身份与维度，并核验了这二百三十六次压力运行保存的匹配见证；该范围内有效见证为二百三十六个，核验失败列表为空。这里的“有效”限定为所检查的匹配见证有效，不表示全部运行均获得最优解，也不表示计时过程已经独立核验。

原作者完整算法流程在已有的七个条件上分别使用十个固定种子，共产生
\begin{equation}
 7\times10=70=65+5
 \label{s0104:author_runs}
\end{equation}
次运行，其中六十五次有效、五次失败。该流程包括至多十轮归约、$S_{\mathrm{pin}}$贪心、以连续未改进次数为停止依据的ILS及展开步骤；ILS的阈值设为$15$，出现改进后计数重新开始，原程序在计数超过阈值时停止，因此不能将其理解为总共只执行十五次迭代。这些配置标识实际测试的方法范围，不与仅运行局部Greedy和ILS的方法混称。七个条件已经包含在初始三十三个条件中，不能再次加到九十二个条件之上。五次失败也属于完整运行结果，不能从报告中删去后把六十五次有效运行写成全部尝试。同样，二十一次重复测时只表示计时观测的次数，并不构成二十一个独立实例。

这些计数反映的是现有材料的覆盖程度。它们没有给出独立数据集总数，也没有证明不同派生条件之间统计独立。分析随机性时，应区分同一输入上的种子变化与不同输入之间的差异；分析运行时间时，还须说明硬件、并发、热启动、停止条件及计时范围。历史实验、压力实验和原作者完整流程的这些设置存在差别，不能把各自记录的时间直接合并成等预算速度排名。早期十秒预算或首次出现正收益的结果仍属于其原有历史口径。

\subsection{从数值输出到可核验结论}

对纯组合模型，一个结果见证至少应包含被选候选，使核验者能够从输入重新计算资源负荷和收益。沿用资源集合$V$与候选集合$E$，令$A_{ve}\in\{0,1\}$表示候选$e\in E$是否占用资源$v\in V$，$b_v\in\mathbb Z_{\geq0}$表示可用份数，$w_e\in\mathbb R_{>0}$表示固定的无量纲收益分值。对输出的二元选择$\widehat x\in\{0,1\}^{|E|}$，定义
\begin{align}
 \ell_v(\widehat x)
 &=\sum_{e\in E}A_{ve}\widehat x_e,
 &&v\in V,
 \label{s0104:checked_load}\\
 W(\widehat x)
 &=\sum_{e\in E}w_e\widehat x_e.
 \label{s0104:checked_value}
\end{align}
其中，$\ell_v$是资源$v$的占用份数，为无量纲非负整数；$W$与$w_e$采用相同的收益分值口径。匹配见证的可行性检查要求对所有$v\in V$均有$\ell_v(\widehat x)\leq b_v$，目标核验则要求重新计算的$W(\widehat x)$与所报告的收益一致。这两项检查分别回答“是否违反资源限制”和“收益是否算对”，都不单独回答“是否还存在更好的选择”。当前组合实验全部采用$b_v=1$，因此相同资源不能被重复占用；这一实验范围不能被扩展解释为一般容量情形的完整验证。

十二候选任务教学例提供了比单个可行见证更充分的最优性依据。其全部二元选择共有$2^{12}=4096$个，独立枚举得到六十二个可行选择，最大收益为$42$，唯一最优选择由\texttt{T02}、\texttt{T03}、\texttt{T04}和\texttt{T12}组成。枚举覆盖了该例的全部二元定义域，所以能够排除收益更高的可行选择。相同输入上的Streaming见证可行，但收益为$39$。这个差异应当保留：它表明该次Streaming输出没有达到该例的最优值，而不能被输入可行性检查所掩盖。上述核验属于确定性见证检查和有限枚举，不是新增计时实验。

较大实例若没有完整枚举，就必须依据另外的界或证明材料讨论最优性。求解器实时报告的目标值、界和停止状态可以描述搜索进展，但实时快照不能替代保存决策后的独立检查。现有算术核验的范围是输入身份、维度和保存的匹配见证，不包含运行时间测量机制或求解器最优性证明。因而，核验记录中没有失败，不能被改写为所有性能指标都已验证。类似地，较大的超边秩描述一个候选涉及较多资源，并不自动证明不可约高阶耦合；当前容量一模型的二元可行域能够用两两冲突表达，算法优势仍须由相应比较结果支持。

\subsection{巡检区域材料的证据边界}

区域巡检材料使用OSM地理参考，并结合目标几何模板和参考物理参数构建场景。当前区域包含$237$座杆塔、$2133$项观测任务、$2$辆车、$4$架无人机和$8$组物理电池对，班次时长为$28800\,\mathrm{s}$；这些数量描述模型配置，不代表已经完成同等规模的现场巡检。真实地理、模板目标和参考设备参数具有不同来源，尚不能合称现场标定数据或试飞结果。现有记录已经包含该区域有限数值转录的构建与运行。较早一次运行出现LP求解器错误，后续两次绑定运行因资源保护退出。后两次运行的构建阶段均对映射并调整地形初值后的四塔暖启动进行了数值检查，记录显示三十六项观测通过了原始代数、索引约束与连续验收；退出时另行保存的候选向量则因资源余量不足而尚未重新验收。两类记录应分别保留，不能仅凭目标值相同就相互继承验收状态。四塔暖启动的通过范围是区域问题中的部分完成方案，既不是完整区域覆盖解，也不是全连续域最优性证明或形式化定向舍入证书。四塔候选的执行见证与覆盖$237$座杆塔的区域模型是不同层次的对象，不能以局部候选的验收代替区域完整覆盖的验收。另一方面，原模型允许部分完成，运行错误及未取得完整覆盖方案也不能推出原联合问题无可行解。实现状态、构建规模、候选可行性、覆盖程度和最优性须分别陈述。
